\documentclass[12pt]{article}

% ======================================================
% Codificación e idioma
% ======================================================

\usepackage[utf8]{inputenc}
\usepackage[T1]{fontenc}
\usepackage{multicol}
\usepackage[spanish,es-nodecimaldot]{babel}

% =====================================================
% Matemáticas
% =====================================================

\usepackage{amsmath,amssymb,amsfonts}

% =====================================================
% Formato y utilidades
% =====================================================

\usepackage[letterpaper,margin=2.2cm]{geometry}
\usepackage{enumitem}
\usepackage{array}
\usepackage{tabularx}
\usepackage{xcolor}
\usepackage{fancyhdr}
\usepackage{lastpage}
\usepackage{hyperref}

% =====================================================
% Hipervínculos
% =====================================================

\hypersetup{
  colorlinks=true,
  linkcolor=blue,
  urlcolor=blue,
  citecolor=blue
}

% =====================================================
% Formato de párrafos
% =====================================================

%\setlength{\parindent}{0pt}
%\setlength{\parskip}{0.5em}

% =====================================================
% Datos editables del examen
% =====================================================

\newcommand{\institucion}{Tecnológico Nacional de México}
\newcommand{\materia}{Cálculo Diferencial}
\newcommand{\nombreexamen}{Segundo Examen Parcial}
\newcommand{\unidad}{Unidad 2: Funciones (pares)}
\newcommand{\docente}{Carlos Ernesto Martínez}
%\newcommand{\fechaexamen}{}
%\newcommand{\duracion}{60 minutos}
%\newcommand{\puntajetotal}{100 puntos}

% =====================================================
% Encabezado y pie de página
% =====================================================

\setlength{\headheight}{15pt}

\pagestyle{fancy}
\fancyhf{}

\fancyhead[L]{\institucion}
\fancyhead[R]{\materia}

\fancyfoot[L]{\docente}
\fancyfoot[C]{Página \thepage\ de \pageref{LastPage}}
\fancyfoot[R]{\nombreexamen}

\renewcommand{\headrulewidth}{0.4pt}
\renewcommand{\footrulewidth}{0.4pt}

% =====================================================
% Contador y ambiente para problemas
% =====================================================

\newcounter{problema}

\newenvironment{problema}[1]
{
  \refstepcounter{problema}
%  \par\medskip
  \noindent
  \textbf{Problema \theproblema.}
  \textbf{[#1 puntos]}
%  \par\smallskip
}
{
%  \par\medskip
}

% =====================================================
% Comando para espacio de respuesta
% =====================================================


% =====================================================
% Documento
% =====================================================

\begin{document}

% =====================================================
% Encabezado principal
% =====================================================

\begin{center}
    {\Large\textbf{\institucion}}\\[0.3em]
    {\large\textbf{\materia}}\\[0.3em]
    {\large\textbf{\nombreexamen}}\\[0.3em]
    {\normalsize\unidad}
\end{center}

%\vspace{0.5em}

% =====================================================
% Datos del estudiante
% =====================================================

\begin{tabularx}{\textwidth}{>{\bfseries}l X >{\bfseries}l X}
Nombre: & \hrulefill &
No. de control: & \hrulefill \\[1.2em]

Grupo: & \hrulefill &
Fecha: & \hrulefill \\[1.2em]

Calificación: & \hrulefill &

\end{tabularx}

%\vspace{1em}

% =====================================================
% Información general
% =====================================================

% =====================================================
% Instrucciones
% =====================================================

\subsection*{Instrucciones}
Lea cuidadosamente cada problema antes de responder. Muestre de manera clara y ordenada todo el procedimiento. Justifique sus resultados cuando se solicite. Las respuestas sin procedimiento no serán consideradas para la calificaci\'on. 
%\medskip


% =====================================================
% REACTIVOS
% =====================================================

\begin{problema}{10}
Determine si el conjunto es una función. Si es una función determine su dominio (Resolver a lo más dos de los tres)
\begin{enumerate}[label=(\alph*)]
\begin{multicols}{3}
\item $\{(x,y)\mid y=\sqrt{x+1}\}$
\item $\{(x,y)\mid y=(x-1)^2+2\}$
\item $\{(x,y)\mid x=(y-2)^2+1\}$

\end{multicols}
\end{enumerate}

\end{problema}

\begin{problema}{15}
Dada $f(x)=2x^2+5x-3$, determine
\begin{enumerate}[label=(\alph*)]
\begin{multicols}{3}
\item $f(h+1)$;
%\item $f(2x^2)$; \item $f(x^2-3)$; \item $f(x+h)$; 
\item $f(x)+f(h)$;
\item $\dfrac{f(x+h)-f(x)}{h}$, $h\ne0$.
 \end{multicols}
\end{enumerate}

\end{problema}

\begin{problema}{15}
Dada $f(x)=\dfrac{3}{x}$, calcule

\begin{enumerate}[label=(\alph*)]
\begin{multicols}{3}
%\item $f(1)$; \item $f(-3)$; \item $f(6)$; \item 
%$f\!\left(\dfrac13\right)$;
\item $f\!\left(\dfrac3a\right)$; %\item $f\!\left(\dfrac3x\right)$;
%\item $\dfrac{f(3)}{f(x)}$; \item $f(x-3)$;
 \item $f(x)-f(3)$;
\item $\dfrac{f(x+h)-f(x)}{h}$, $h\ne0$.
\end{multicols}
\end{enumerate}
\end{problema}

\begin{problema}{20}
Dibuje a mano la gráfica de la función y determine su dominio y su imagen (Resolver a lo más dos de los tres).

\begin{enumerate}[resume]
\begin{multicols}{3}
%\item $f(x)=3x-1$
%\item $g(x)=4-x$
%\item $F(x)=2x^2$
%\item $G(x)=x^2+2$
%\item $g(x)=5-x^2$
\item $f(x)=(x-1)^2$
%\item $G(x)=\sqrt{x-1}$
%\item $F(x)=\sqrt{9-x}$
%\item $f(x)=\sqrt{x^2-4}$
%\item $g(x)=\sqrt{4-x^2}$
%\item $g(x)=\sqrt{9-x^2}$
%\item $f(x)=\sqrt{x^2-1}$
\item $G(x)=\begin{cases}9-x^2,&x\ne-3,\\4,&x=-3.\end{cases}$
\item $F(x)=\begin{cases}x^2-4,&x<3,\\2x-1,&3\le x.\end{cases}$
\end{multicols}
\end{enumerate}
\end{problema}

\begin{problema}{25}
Dibuje la gráfica de la función escalón (o salto) unitario denotada por $U$ y definida por
\[
U(x)=\begin{cases}0,&x<0,\\1,&0\le x.\end{cases}
\]
Defina cada una de las siguientes funciones a trozos y dibuje sus gráficas:
\begin{enumerate}
\begin{multicols}{3}
\item $U(x-1)$;  \item $U(x)-U(x-1)$. \item $(x+1)U(x+1)$; 
\end{multicols}
\end{enumerate}
\end{problema}

\begin{problema}{25}
Dibuje la gráfica de la función signo denotada por $\operatorname{sgn}$ y definida por
\[
\operatorname{sgn}x=\begin{cases}-1,&x<0,\\0,&x=0,\\1,&0<x.\end{cases}
\]
$\operatorname{sgn}(x)$ se lee ``signo de $x$''. Defina cada una de las siguientes funciones a trozos y dibuje sus gráficas:
\begin{enumerate}
\begin{multicols}{3}
\item $x\,\operatorname{sgn}(x)$; \item $x-2\,\operatorname{sgn}(x)$.
\item $\operatorname{sgn}(x+1)-\operatorname{sgn}(x-1)$.
\end{multicols}
\end{enumerate}

\end{problema}

\begin{problema}{20}
defina las siguientes funciones y determine el dominio de la función resultante: (a) $f+g$; (b) $f-g$; (c) $f\cdot g$; (d) $f/g$; (e) $g/f$. (Resolver a lo más una de las dos)

\begin{enumerate}
\begin{multicols}{2}
\item $f(x)=\sqrt{x}$; $g(x)=4-x^2$
\item $f(x)=x^2$; $g(x)=\dfrac1{\sqrt{x}}$\end{multicols}
\end{enumerate}

\end{problema}

\begin{problema}{20}
Defina las siguientes funciones y determine el dominio de la función compuesta: (a) $f\circ g$; (b) $g\circ f$; (c) $f\circ f$; (d) $g\circ g$. (Resolver a lo más uno de los dos)

\begin{enumerate}
\begin{multicols}{2}
\item $f(x)=3-2x$; $g(x)=6-3x$
%\item $f(x)=x^2-1$; $g(x)=\dfrac1x$
\item $f(x)=\sqrt{x^2-1}$; $g(x)=\sqrt{x-1}$
\end{multicols}

\end{enumerate}

\end{problema}

\begin{problema}{15}
Si $f$ y $g$ son dos funciones tales que $(f\circ g)(x)=x$ y $(g\circ f)(x)=x$, entonces se dice que $f$ y $g$ son inversas una de la otra. Demuestre que $f$ y $g$ son inversas una de la otra. (Resolver a lo más uno de los dos)

\begin{enumerate}
\begin{multicols}{2}

\item $f(x)=\dfrac1{x+1}$ y $g(x)=\dfrac{1-x}{x}$
\item $f(x)=x^2$, $x\le0$, y $g(x)=-\sqrt{x}$

\end{multicols}
\end{enumerate}
\end{problema}

\begin{problema}{15}

Determine si cada función es inyectiva en el dominio indicado. (Resolver a lo más uno de los dos) Justifique.
\begin{enumerate}
\begin{multicols}{2}

\item $f:\mathbb R\to\mathbb R$, $f(x)=x^2$.

%\item $f:[0,\infty)\to\mathbb R$, $f(x)=x^2$.

\item $f:[0,\infty)\to[0,\infty)$, $f(x)=\sqrt{x}$.


\end{multicols}
\end{enumerate}

\end{problema}

\begin{problema}{15}
Encuentre la inversa, especificando las restricciones necesarias. (Resolver a lo más dos de los tres)
\begin{enumerate}
\begin{multicols}{3}

\item $f(x)=\dfrac{x+1}{x-2}$.

\item $f(x)=\dfrac{3x+1}{2x-5}$.

\item $f(x)=1+\dfrac1x$.
\end{multicols}
\end{enumerate}

\end{problema}

\begin{problema}{15}
Encuentre $f^{-1}$ y compruebe $f(f^{-1}(x))=x$ y $f^{-1}(f(x))=x$. (Resolver a lo más dos de los tres)
\begin{enumerate}
\begin{multicols}{3}
\item $f(x)=\dfrac{x-5}{2}$.
\item $f(x)=x^3+2$.
\item $f(x)=e^{2x}$.
\end{multicols}
\end{enumerate}

\end{problema}


% =====================================================
% Tabla de calificación
% =====================================================
\begin{center}
\textbf{Calificación final:}\qquad
\rule{3cm}{0.4pt}
\end{center}


\end{document}